% Lecture example set: SC4061-L02
% Source: L1 pp. 27, 31--34 and L1_MATLAB.m; E02--E03 are self-contained checks
% Human verified: Yes

\exercisemeta
  {SC4061-L02-EX}
  {2026-08-18}
  {L1 Imaging pp.~27, 31--34；L1\_MATLAB.m；课堂讨论中的自编检验例}
  {答案已完成并确认}
  {已确认（2026-08-18）}

\exercisequestion{SC4061-L02-E01}{Perceptron 权重更新}

\textbf{类型：}lecture example（讲义例题）。

\textbf{对应知识点：}
\relatedtopic{SC4061-L02-T03}{Perceptron Learning Rule}

\subsubsection*{题目}

给定增广样本
\[
  \tilde{\bm x}_1=[3,3,1]^{\mathsf T}\in C_1,
  \qquad
  \tilde{\bm x}_2=[1,1,1]^{\mathsf T}\in C_2,
\]
依次循环输入 $\tilde{\bm x}_1,\tilde{\bm x}_2$。令 $\alpha=1$、$\tilde{\bm w}(1)=\bm 0$，使用本讲 Perceptron learning rule 求 $\tilde{\bm w}(4)$、$\tilde{\bm w}(5)$ 及最终 decision boundary（决策边界）。

\begin{checkedsolution}
前两次更新为
\[
  \tilde{\bm w}(2)=\bm0+\tilde{\bm x}_1=[3,3,1]^{\mathsf T},
\]
\[
  \tilde{\bm w}(3)=\tilde{\bm w}(2)-\tilde{\bm x}_2
  =[2,2,0]^{\mathsf T}.
\]
当 $k=3$ 时，$\tilde{\bm w}^{\mathsf T}(3)\tilde{\bm x}_1=12>0$，分类正确，因此
\[
  \boxed{\tilde{\bm w}(4)=[2,2,0]^{\mathsf T}}.
\]
当 $k=4$ 时，$\tilde{\bm w}^{\mathsf T}(4)\tilde{\bm x}_2=4\ge0$，$C_2$ 样本被误分，因此
\[
  \boxed{\tilde{\bm w}(5)=[2,2,0]^{\mathsf T}-[1,1,1]^{\mathsf T}
  =[1,1,-1]^{\mathsf T}}.
\]
继续循环后权重稳定为 $[1,1,-3]^{\mathsf T}$，故
\[
  \boxed{x_1+x_2-3=0}.
\]
检验：$(3,3)$ 的 score 为 $3>0$，属于 $C_1$；$(1,1)$ 的 score 为 $-1<0$，属于 $C_2$。
\end{checkedsolution}

\textbf{检查点：}第三个分量是 bias weight（偏置权重），不能在更新时遗漏增广向量最后的 $1$。

\exercisequestion{SC4061-L02-E02}{LiDAR Time-of-Flight 测距}

\textbf{类型：}self-contained check example（自编检验例，基于讲义 LiDAR 场景）。

\textbf{对应知识点：}
\relatedtopic{SC4061-L02-T01}{Beyond a Conventional 2D Camera}

\subsubsection*{题目}

LiDAR 发射脉冲后经过 $\Delta t=20\,\mathrm{ns}$ 收到回波。取 $c=3.0\times10^8\,\mathrm{m/s}$，估计物体距离。

\begin{checkedsolution}
\[
  d=\frac{c\Delta t}{2}
  =\frac{(3.0\times10^8)(20\times10^{-9})}{2}
  =\boxed{3\,\mathrm m}.
\]
\end{checkedsolution}

\textbf{检查点：}$c\Delta t=6\,\mathrm m$ 是 pulse（脉冲）的往返路程，不是传感器到物体的单程距离。

\exercisequestion{SC4061-L02-E03}{RGB Threshold 与 Logical Mask}

\textbf{类型：}self-contained check example（自编检验例，基于 L1\_MATLAB.m）。

\textbf{对应知识点：}
\relatedtopic{SC4061-L02-T06}{Logical Mask 与 Vectorized Operation}

\subsubsection*{题目}

规则把满足 $R+G+B>300$ 的 pixel（像素）改为 $(60,60,60)$，否则保持不变。求
\[
  p_1=(120,110,90),
  \qquad
  p_2=(100,90,80)
\]
处理后的值。

\begin{checkedsolution}
$p_1$ 的分量和为 $320>300$，mask value（掩膜值）为 $1$，因此
\[
  \boxed{p_1'=(60,60,60)}.
\]
$p_2$ 的分量和为 $270\le300$，mask value 为 $0$，因此保持不变：
\[
  \boxed{p_2'=(100,90,80)}.
\]
\end{checkedsolution}

\textbf{检查点：}判断条件是严格的大于号；若分量和恰好为 $300$，该 pixel 不会被替换。
